\documentclass[conference,onecolumn]{IEEEtran}

\usepackage[T1]{fontenc}
\usepackage[utf8]{inputenc}
\usepackage[spanish,es-tabla]{babel}
\usepackage{lmodern}
\usepackage{microtype}
\usepackage{csquotes}
\usepackage{amsmath,amssymb}
\usepackage{booktabs}
\usepackage{float}
\usepackage{tabularx}
\usepackage{array}
\usepackage{longtable}
\usepackage{url}
\usepackage[hidelinks]{hyperref}
\usepackage[backend=biber,style=ieee,sorting=none]{biblatex}
\addbibresource{referencias.bib}
\AtBeginDocument{\spanishplainpercent}

\newcolumntype{Y}{>{\raggedright\arraybackslash}X}
\newcolumntype{L}[1]{>{\raggedright\arraybackslash}p{#1}}
\renewcommand{\IEEEkeywordsname}{Palabras clave}

\title{FORM-TL-01: Taller de tipos de requisitos para la Plataforma de gestión de citas médicas del Policlínico Ribera del Quevedo}
\author{\IEEEauthorblockN{Danela Dayana Arteaga Alava}\IEEEauthorblockA{Carrera de Software\\
Facultad de Ciencias de la Computación\\
Universidad Técnica Estatal de Quevedo, Ecuador\\
Ingeniería de Requerimientos -- ISR-401}}

\begin{document}
\maketitle

\begin{abstract}
Actividad formativa individual de la Semana 1. Se clasifican los 20 enunciados del caso «Plataforma de gestión de citas médicas» en tres dimensiones: tipo de requisito (funcional, de calidad o restricción), producto o proceso, y sistema o software. A los requisitos de calidad se les asigna su característica y subcaracterística según ISO/IEC 25010:2023.
\end{abstract}

\begin{IEEEkeywords}
Ingeniería de requisitos, tipos de requisitos, ISO/IEC 25010:2023, requisitos de calidad.
\end{IEEEkeywords}

\section*{Datos generales}

\begin{table}[H]
\centering\footnotesize
\begin{tabularx}{\textwidth}{@{}p{0.22\textwidth}Y@{}}
\toprule
\textbf{Campo} & \textbf{Detalle}\\
\midrule
Asignatura & Ingeniería de Requerimientos (ISR-401) · 4.\textsuperscript{o} nivel · Paralelo B\\
Estudiante & Arteaga Alava Danela Dayana\\
Docente & Gleiston Cicerón Guerrero Ulloa, PhD\\
Fecha de cierre & Sábado 7 de noviembre de 2026, 23:59\\
Repositorio & \url{https://github.com/darteagaa-boop/ISR401-ArteagaDanela}\\
\bottomrule
\end{tabularx}
\end{table}

\section{Criterios de clasificación}
Cada enunciado se clasifica en tres dimensiones, siguiendo la distinción usada en la ingeniería de requisitos \cite{ieee29148,glinz2007,sommerville2016}:

\begin{itemize}
\item \textbf{Tipo:} funcional, que se refiere a lo que hace el sistema; de calidad, que hace énfasis en qué tan bien lo hace, y en este caso se asigna su característica y subcaracterística de ISO/IEC 25010:2023 \cite{iso25010}; o restricción, que es un límite impuesto a la solución o al proyecto.
\item \textbf{Producto o proceso:} se refiere al sistema que se construye o a la forma de desarrollarlo, ya sea en plazos, validaciones o revisiones.
\item \textbf{Sistema o software:} es de sistema si para cumplirse intervienen elementos además del software, como el hardware, servicios externos, infraestructura o procedimientos de la organización; en caso contrario, es de software.
\end{itemize}

En el modelo de calidad de producto de ISO/IEC 25010:2023 las características son: adecuación funcional, eficiencia de desempeño, compatibilidad, capacidad de interacción, fiabilidad, seguridad, mantenibilidad, flexibilidad y protección \cite{iso25010}.

\clearpage
\section{Tabla de clasificación de los 20 enunciados}
La Tabla~\ref{tab:clasificacion} presenta la clasificación de cada enunciado con su justificación.

{\footnotesize
\setlength{\tabcolsep}{3pt}
\begin{longtable}{@{}L{0.04\textwidth}L{0.18\textwidth}L{0.11\textwidth}L{0.15\textwidth}L{0.08\textwidth}L{0.08\textwidth}L{0.27\textwidth}@{}}
\caption{Clasificación de los enunciados de la Plataforma de gestión de citas médicas}\label{tab:clasificacion}\\
\toprule
\textbf{ID} & \textbf{Enunciado} & \textbf{Tipo} & \textbf{Característica; subcaracterística ISO/IEC 25010:2023} & \textbf{Producto / Proceso} & \textbf{Sistema / Software} & \textbf{Justificación}\\
\midrule
\endfirsthead
\toprule
\textbf{ID} & \textbf{Enunciado} & \textbf{Tipo} & \textbf{Característica; subcaracterística ISO/IEC 25010:2023} & \textbf{Producto / Proceso} & \textbf{Sistema / Software} & \textbf{Justificación}\\
\midrule
\endhead
\midrule
\multicolumn{7}{r@{}}{\emph{Continúa en la página siguiente}}\\
\endfoot
\bottomrule
\endlastfoot
E01 & Reservar cita eligiendo especialidad, médico, fecha y hora. & Funcional & -- & Producto & Software & Describe un servicio que el software ofrece al paciente.\\
E02 & Recordatorio por SMS 24~h antes de la cita. & Funcional & -- & Producto & Sistema & Es un comportamiento del producto, pero se cumple con un servicio externo de mensajería, no solo con el software.\\
E03 & La aplicación debe ser rápida. & Calidad (ambiguo) & Eficiencia de desempeño; comportamiento temporal & Producto & Software & Expresa cómo responde el producto, no qué hace; no es verificable tal como está y requiere cuantificación.\\
E04 & El recepcionista registra la llegada del paciente en ventanilla. & Funcional & -- & Producto & Software & Define una acción concreta que el software permite a un actor: registrar la llegada.\\
E05 & Desarrollo en Java 21 con PostgreSQL. & Restricción (tecnológica) & -- & Producto & Software & Limita las opciones de diseño e implementación del producto, pero no es una función ni un atributo de calidad.\\
E06 & El médico consulta la lista del día y marca consultas atendidas. & Funcional & -- & Producto & Software & Describe funciones observables del software.\\
E07 & Los datos de salud deben estar seguros. & Calidad (ambiguo) & Seguridad; confidencialidad e integridad & Producto & Software & Atributo de protección de la información; «seguros» no es medible y requiere cuantificación.\\
E08 & La app móvil funciona en Android 10 o superior. & Restricción (plataforma) & -- & Producto & Software & Fija un entorno operativo mínimo y restringe la solución en lugar de describir una función.\\
E09 & Cita cancelada con menos de 4~h de anticipación = inasistencia. & Funcional (regla de dominio) & -- & Producto & Software & Regla de negocio que el software debe aplicar al procesar cancelaciones.\\
E10 & La plataforma debe ser fácil de usar para adultos mayores. & Calidad (ambiguo) & Capacidad de interacción; aprendibilidad, operabilidad, inclusividad & Producto & Software & Atributo de interacción con un grupo de usuarios; el término «fácil» es subjetivo y requiere cuantificación.\\
E11 & Entrega en un plazo máximo de seis meses. & Restricción (proyecto) & -- & Proceso & Sistema & Restringe el proyecto completo mediante el tiempo.\\
E12 & Cierre automático de sesión tras 10~min sin actividad. & Funcional (derivado de seguridad) & Seguridad; confidencialidad & Producto & Software & Es un comportamiento concreto y ya verificable que nace de una preocupación de seguridad, pero se expresa como función.\\
E13 & Quiosco lee la cédula e imprime ticket con número de turno. & Funcional & -- & Producto & Sistema & Función que involucra hardware (pantalla táctil e impresora térmica) además del software.\\
E14 & La plataforma debe estar siempre disponible. & Calidad (ambiguo) & Fiabilidad; disponibilidad & Producto & Sistema & La disponibilidad depende de servidores, red e infraestructura; «siempre» no es alcanzable ni medible y requiere cuantificación.\\
E15 & El director genera un reporte mensual de inasistencias por especialidad. & Funcional & -- & Producto & Software & Describe una salida concreta que el software produce para un actor.\\
E16 & Cada incremento se valida con el director cada dos semanas. & Restricción (proceso) & -- & Proceso & Sistema & Regula cómo se desarrolla y valida el proyecto, no el producto.\\
E17 & Cumplir la Ley Orgánica de Protección de Datos Personales. & Restricción (legal) & -- & Producto & Sistema & Obligación normativa externa que afecta al software y a los procedimientos de la organización.\\
E18 & Integrarse sin problemas con el sistema de facturación. & Calidad (ambiguo) & Compatibilidad; interoperabilidad & Producto & Sistema & Intercambio de información con otro sistema existente; el término «sin problemas» no es verificable y requiere cuantificación.\\
E19 & Si el servidor falla, restablecer el servicio en menos de 30~min sin perder citas confirmadas. & Calidad (cuantificado) & Fiabilidad; recuperabilidad & Producto & Sistema & Tiene métrica de tiempo de recuperación, unidad y umbral (menos de 30~min); involucra infraestructura.\\
E20 & Todo cambio de código se revisa por otro integrante antes de integrarse. & Restricción (proceso) & -- & Proceso & Software & Prescribe una práctica de desarrollo sobre el código fuente, no un atributo del producto.\\
\end{longtable}
}

\textbf{Resumen:} 8 funcionales (E01, E02, E04, E06, E09, E12, E13, E15), 6 de calidad (E03, E07, E10, E14, E18, E19; los cinco primeros son ambiguos) y 6 restricciones (E05, E08, E11, E16, E17, E20). Son de proceso E11, E16 y E20; el resto es de producto.

\printbibliography[title={Referencias}]
\end{document}
